% Source: https://artofproblemsolving.com/community/c3h1696911_geometry_problemsquares_and_triangles
% Stable identity: aops-topic-1696911-q01
Rotate the figure about $B$ through $90^\circ$ so that $C$ goes to $A$ (this is possible since $BC=BA$ and $\angle CBA=90^\circ$), and let $P$ go to $P'$. Then triangle $BPC$ goes to triangle $BP'A$, so
\[
BP'=BP=\sqrt2,\qquad AP'=CP=1,\qquad \angle BP'A=\angle BPC,\qquad \angle PBP'=90^\circ .
\]

\begin{asy}
usepackage("amsmath");
size(6cm);
defaultpen(fontsize(10pt));
real s=sqrt(5);
pair A=(0,0),B=(s,0),C=(s,s),D=(0,s);
pair P=(4/s,3/s);
pair Pp=B+rotate(90)*(P-B);
draw(A--B--C--D--cycle,linewidth(0.9));
draw(P--A,linewidth(0.7));draw(P--B,linewidth(0.7));draw(P--C,linewidth(0.7));
draw(B--Pp,blue+dashed);draw(A--Pp,blue+dashed);draw(P--Pp,blue+dashed);
pair u=unit(A-Pp),v=unit(P-Pp);real q=0.13;
draw(Pp+q*u--Pp+q*(u+v)--Pp+q*v,blue);
pair u2=unit(P-B),v2=unit(Pp-B);
draw(B+q*u2--B+q*(u2+v2)--B+q*v2,blue);
dot(P);dot(Pp,blue);
label("$A$",A,NW);label("$B$",B,E);label("$C$",C,NE);label("$D$",D,NW);
label("$P$",P,dir(115));label("$P'$",Pp,S,blue);
label("$\sqrt5$",(A+P)/2,NW);label("$\sqrt2$",(B+P)/2,W);label("$1$",(C+P)/2,NW);
label("$1$",(A+Pp)/2,SW,blue);label("$2$",(P+Pp)/2,2*E,blue);
label("\mbox{}",(min(currentpicture,user=true)+max(currentpicture,user=true))/2,black); // keeps the text after the figure black when LaTeX typesets the labels inline
\end{asy}

\textbf{(a)} Triangle $PBP'$ is a right isosceles triangle with legs $\sqrt2$, so $PP'=2$ and $\angle BP'P=45^\circ$. In triangle $AP'P$, $AP'^2+PP'^2=1+4=5=AP^2$, so by the converse of Pythagoras' theorem $\angle AP'P=90^\circ$. The ray $P'P$ lies between $P'B$ and $P'A$ (as in the figure), so
\[
\angle BPC=\angle BP'A=\angle BP'P+\angle PP'A=45^\circ+90^\circ=\boxed{135^\circ}.
\]

\textbf{(b)} By the law of cosines in triangle $BPC$,
\[
BC^2=BP^2+CP^2-2\cdot BP\cdot CP\cos135^\circ=2+1+2\sqrt2\cdot\dfrac{\sqrt2}{2}=5,
\]
so the side length is $\boxed{\sqrt5}$.

(Check with coordinates: $A=(0,0)$, $B=(\sqrt5,0)$, $C=(\sqrt5,\sqrt5)$ and $P=\left(\dfrac{4}{\sqrt5},\dfrac{3}{\sqrt5}\right)$ give $AP^2=5$, $BP^2=\dfrac{1}{5}+\dfrac{9}{5}=2$, $CP^2=\dfrac{1}{5}+\dfrac{4}{5}=1$.)
